\bigskip
\noindent Autres pistes de l'onglet Findings pour ce thème, dans des cotes qui n'ont pas de section ici :\par

\begin{encadre}
\begin{definition*}[Grothendieck, pages 2 à 16, catégorie de Teichmüller--Galois\piste{142-teichmuller-galois-category-etale}]
La catégorie de Teichmüller--Galois a pour objets les courbes anabéliennes de tous types $(g, \nu)$ et pour flèches les complétés profinis des classes d'isotopie d'applications étales finies ; elle est obtenue par une complétion profinie des catégories dont les groupes d'automorphismes vérifient la condition de finitude $(\ast)$, l'image du stabilisateur de chaque flèche $u : X \to Y$ étant d'indice fini dans $\mathrm{Aut}(Y)$. Chaque groupe d'automorphismes y est le groupe de Teichmüller profini ; la construction est algébrique, invariante par extension de corps algébriquement clos de caractéristique~$0$, et $\mathrm{Aut}(k)$ y opère.
\end{definition*}
\end{encadre}
{\small Les flèches de cette catégorie relient des types $(g, \nu)$ différents par des applications étales finies, et non par les inclusions de sous-surfaces et les recollements de la tour de Teichmüller, dont Hatcher, Lochak et Schneps (2000) \cite{hatcherlochakschneps2000} donnent une présentation par les niveaux~$1$ et~$2$ ; l'action de $\mathrm{Aut}(k)$ y est l'analogue, pour tous les types, de l'action de Galois sur $\hat\pi_1(M_{0,4}) = \hat F_2$ de la section~\ref{sec:145}. Mathlib n'a pas la complétion profinie des catégories ni les groupes de Teichmüller. On se reportera aux notes de Grothendieck (\lot{142}{1}{2}).\par}

\begin{encadre}
\begin{theoreme*}[Grothendieck, pages 167 et 175 à 180, monodromie modérée à plusieurs paramètres\piste{34-multiparameter-tame-monodromy}]\nonverifie
Soient $S$ un schéma strictement local de dimension quelconque et $\bar X$ un $S$-schéma projectif et plat, de fibre fermée à croisements normaux hors de la codimension~$2$, lisse à bord à croisements normaux relatifs au-dessus d'un ouvert $U \subset S$. L'image de $\pi_1(U)$ dans $\mathrm{Out}(\pi_1^{(p')})$ d'une fibre géométrique est un groupe abélien d'ordre premier à $p$, de rang au plus le nombre $\nu$ de points doubles en dimension relative~$1$.
\end{theoreme*}
\end{encadre}
{\small Les feuillets en esquissent une démonstration incomplète (pages 175 à 180), où le lemme~3.1, sur la courbe nodale verselle, n'est pas démontré. Les hypothèses (a) à (d) reposent sur des lectures incertaines de la page~167. Sur une base de dimension~$1$, la finitude de l'image de l'inertie sauvage dans $\mathrm{Out}(\pi_1^{(p')})$ d'une fibre géométrique, pour toute variété sur un corps complet de valuation discrète, est le théorème~2.1 de Kisin (2000) \cite{kisin2000}, démontré par les altérations et non par réduction aux courbes. En dimension relative~$1$, la lecture y voit la description de la monodromie le long d'une courbe stable dégénérée par les twists de Dehn des points doubles, élevés aux multiplicités~$n_i$ des équations locales $xy = t^{n_i}$ ; en dimension relative $\geq 2$, l'énoncé encadré reste ouvert ici. Quand $S \smallsetminus U$ est un diviseur à $r$ branches à croisements normaux, le groupe fondamental modéré de $U$ au voisinage du point fermé est $\prod_{\ell \neq p} \mathbf Z_\ell(1)^r$ (SGA~1, exposé~XIII, appendice~I, corollaire~5.3 \cite{sga1}), abélien d'ordre premier à~$p$ ; la section~\ref{sec:29} en calcule la fonctorialité. Mathlib n'a pas la monodromie modérée. On se reportera aux notes de Grothendieck (\lot{34}{9}{167}).\par}

\begin{encadre}
\begin{contrexemple*}[pages 36 à 39, le lemme du Complément\piste{149-dominant-faithfulness-curves}]\nonverifie
Le lemme de la page~36, qui omet l'hypothèse de~$(\ast)$ sur les groupes fondamentaux, est faux pour les courbes hyperboliques de types $(1,1)$ et $(1,2)$, et réduire l'ouvert $V$ n'écarte pas le cas exceptionnel $\operatorname{card} I = 1$. L'extension que la page annonce aux polycourbes hyperboliques de toute dimension repose sur un plongement dans une variété semi-abélienne, qui manque en dimension~$2$ : la courbe elliptique universelle privée de sa section nulle, au-dessus de $Y(N)$, n'en a pas.
\end{contrexemple*}
\end{encadre}
{\small Pour les courbes, l'énoncé de la piste, en caractéristique~$0$ sur un corps algébriquement clos (deux morphismes $f, g : X \to Y$ vers une courbe hyperbolique, $f$ dominant, qui induisent le même homomorphisme extérieur des groupes fondamentaux étales profinis, sont égaux), n'est pas démontré ici ; l'édition ferme les types $(1,1)$ et $(1,2)$ par l'hypothèse de~$(\ast)$. Mathlib n'a pas le groupe fondamental étale des courbes hyperboliques. On se reportera aux notes de Grothendieck (\lot{149}{2}{21}).\par}

\begin{encadre}
\begin{theoreme*}[Grothendieck, pages 153 à 155, 1-asphéricité locale\piste{34-lci-local-1-asphericity}]\nonverifie
Soit $f : X \to S$ plat sur un trait strictement local, de fibre fermée $X_0$ localement d'intersection complète, dont le lieu non lisse est de codimension $\geq 2$ dans $X_0$. Alors $f$ est localement 1-asphérique pour les premiers $\neq p$ ; si $f$ est propre, la spécialisation $H^1(X_0, G) \to H^1(X_{\bar\eta}, G)$ est bijective pour tout groupe fini $G$ d'ordre premier à~$p$.
\end{theoreme*}
\end{encadre}
\begin{demopage}[pages 153 à 155]
Au n\textsuperscript{o}~5.5, $f$ étant localement 1-acyclique aux générisations strictes de $x$, la 1-asphéricité locale en $x$ pour un groupe fini $G$ d'ordre premier à $p$ se ramène à (5.5.1), $H^i(\bar V, G) \simeq H^i(V_{\bar\eta}, G)$ pour $i \leq 1$ (p.~153), puis à (5.5.2), $H^i(V(h), G) = 0$ pour $i \leq 1$ (p.~154). La page affirme sans démonstration (\q{On sait}) que (5.5.2) vaut aux points d'intersection complète de dimension $\geq 3$, et la lecture le rapporte à la pureté de SGA~2. Pour $f$ plat, cela revient à $X_0$ d'intersection complète de dimension $\geq 2$ en $x$, et le cas lisse traite les autres points, d'où le théorème~5.6 et, sans autre argument, l'assertion sur $H^1$.
\end{demopage}

\begin{encadre}
\begin{conjecture*}[D'après les pages 94, 95 et 98 à 101, forme regroupée par l'édition, image de Galois sur une seule courbe\piste{142-galois-image-conditions-single-curve}]
Soient $X$ une courbe algébrique anabélienne sur un corps $K$ de type fini sur $\mathbf Q$ et $\Theta$ un automorphisme du groupoïde fondamental profini $\hat\Pi_1(X, I)$ qui fixe les points base. Si $\Theta$ se relève à $\hat\Pi_1(X', I')$ pour tout revêtement étale fini $X'$ de $X$ défini sur $K$, respecte les noyaux de $\hat\pi_1(X', x') \to \hat\pi_1(Y, f(x'))$ pour tout $K$-morphisme $f : X' \to Y$, est compatible aux morphismes entre revêtements et vérifie la condition des lacets avec un seul multiplicateur, alors $\Theta$ provient de $\mathrm{Gal}(\bar K/K)$. Ces conditions sont nécessaires.
\end{conjecture*}
\end{encadre}
{\small La page~95 conjecture suffisantes les conditions de relèvement et de noyaux ; leur regroupement avec les conditions des pages 98 à 101 est de l'édition. En termes actuels, c'est une forme de la question d'Ihara et de la conjecture d'Oda--Matsumoto (I/OM) attachée à une seule courbe et à ses revêtements : I/OM demande si $\mathrm{Gal}(\overline K/K)$ est le groupe des automorphismes du foncteur groupe fondamental géométrique sur une catégorie de $K$-variétés. Topaz (2018) \cite{topaz2018} en démontre une forme renforcée, pour le quotient modulo~$\ell$ abélien-par-central de ce groupe, sur un corps parfait infini de caractéristique $\neq \ell$ et pour toute catégorie de variétés assez grosse au sens de son §~1.7 ; cette forme porte sur une catégorie de variétés, non sur une courbe et ses revêtements, et ne fait pas intervenir de condition d'inertie. Les variantes tripodes finies de Pop (2019) \cite{pop2019} supposent une base de dimension au moins deux, munie d'un nombre fini d'applications rationnelles dominantes vers $\mathbf P^1 \smallsetminus \{0,1,\infty\}$ ; elles ne couvrent donc pas les revêtements d'une seule courbe, et c'est la condition des lacets de la page qui devrait en tenir lieu. Pour $X = \mathbf P^1 \smallsetminus \{0, 1, \infty\}$, la condition des lacets est celle des sous-groupes $L_0$, $L_1$, $L_\infty$ de la section~\ref{sec:145}. Mathlib n'a pas les groupoïdes fondamentaux profinis des courbes. On se reportera aux notes de Grothendieck (\lot{142}{5}{94}).\par}

\begin{encadre}
\begin{conjecture*}[Grothendieck, pages 72 à 77 et 84 à 93, spécialisation vers une pointe\piste{149-cusp-specialisation-index}]
Soient $U = X \setminus T$ une courbe hyperbolique relative sur un trait et $f$ une section rationnelle qui rencontre la section de pointe $g_{i_0}$ dans la fibre spéciale avec la multiplicité $n$. L'inertie s'envoie sur $n$ fois le sous-groupe de lacet $L_{i_0}$ de cette pointe, et elle agit trivialement exactement quand $f$ se prolonge.
\end{conjecture*}
\end{encadre}
{\small Mathlib n'a pas la spécialisation des sections des courbes hyperboliques relatives. On se reportera aux notes de Grothendieck (\lot{149}{4}{72}).\par}
